\newcommand{\ThreeTensor}[1]{\mathsf{#1}}

The ramification of the rigid cross section assumption is that the normal and shear strain components in the plane of a cross section vanish; in our coordinate system that means
\(E_{yy}=E_{zz}=0\) and \(E_{yz}=E_{zy}=0\). 
For isotropy, this
constraint is not a problem in the shear strain components because the
shear stresses and strains are uncoupled under isotropy (for example,
\(S_{yz}=2 \mu E_{yz}\) ). Therefore, vanishing shear strains simply
implies vanishing shear stress.

The constitutive equations for isotropic hyperelasticity have an
inherent coupling of the normal strains and stresses. The stresses are
given by \[
\mathbf{S}=\lambda \operatorname{tr}(\bm{E}) \bm{1}+2 \mu \mathbf{E}
\] Since under the strain hypothesis
\(\operatorname{tr}(\bm{E})=E_{xx}\), we have
\(S_{yy}=S_{zz}=\lambda E_{xx}\). Now, from the constitutive law
\(S_{xx}=(\lambda+2 \mu) E_{xx}\) so \(S_{yy}=S_{zz}=\nu S_{xx}\), where
\(\nu\) is Poisson's ratio. Thus, the constraint induces normal stresses
in the plane of the cross section owing to Poisson's effect (tension in
one direction causes lateral contraction of the dimensions perpendicular
to the direction of tension). Observational evidence on the behavior of
beams would indicate that these stresses tend to be rather small. In
fact, it is possible that the presence of these stresses will violate
the traction boundary conditions on the lateral surface of the beam. For
a beam with a traction-free lateral surface, we can argue that the
normal stresses \(S_{yy}\) and \(S_{zz}\) are very small because the
cross-sectional dimensions are small compared with the length of the
beam. We would like to make the assumption that \(S_{yy}=S_{zz}=0\).

Let us examine what would have happened if we had made the assumption of
vanishing normal stress and not the assumption of vanishing normal
strain. We made exactly that assumption for the uniaxial tension test in
order to recast the constitutive equations in terms of Young's modulus
and Poisson's ratio. When we made the assumption that
\(S_{yy}=S_{zz}=0\), we got the constitutive relationship
\(S_{xx}=C E_{xx}\), rather than \(S_{xx}=(\lambda+2 \mu) E_{xx}\),
which results for the vanishing strain assumption. This gives us a way
to partially recover from our difficulties. In the constitutive
relationships for beams, the quantity \(E=\lambda+2 \mu\) appears
repeatedly. If we simply substitute the value of Young's modulus \(E=C\)
instead, then the results of beam theory accord well with observation.

\subsection{Stress Space}\label{stress-space}

Consider \(\mathbb{S}_{\mathrm{rod}}\subset \mathbb{S}\) 
\[
\operatorname{dim}\mathbb{S}_{\mathrm{rod}} = 6 - 3 = 3
\]

\subsubsection{Coordinates}\label{coordinates}

Let \(\bm{\sigma} \triangleq \bm{S}\FrameBasis\) and define
\(\bm{\varepsilon}\) so that
\(\FourTensor{P}_{\mathrm{rod}}\bm{E} = \FourTensor{P}_{\mathrm{sym}}[\FrameBasis \otimes\bm{\varepsilon}]\).
Note that the shear components of \(\bm{\varepsilon}\) and \(\bm{E}\)
are related by 
\[
\FrameBasis_{\beta}\cdot\bm{\varepsilon} = \FrameBasis_{\beta} \cdot 2\bm{E}\FrameBasis
\] 
which reads, in matrix notation,
\((\varepsilon_2, \varepsilon_3) = (2E_{21},\, 2E_{31})\).

The components \(\sigma_k\) of \(\bm{\sigma}\) on the frame basis
\(\{\FrameBasis_k\}\) are equivalent to the components of
\(\bm{S}\in\mathbb{S}_{\mathrm{rod}}\) on a basis \(\{\bm{I}_k\}\)
defined by \[
\bm{I}_k \triangleq \FrameBasis\otimes\FrameBasis_k
\]
%
An alternative basis for \(\mathbb{S}_{\mathrm{rod}}\) is furnished by
the three tensors 
\[
\hat{\bm{I}} = \FrameBasis \otimes \FrameBasis,\, 
\quad\text{ and }\quad 
\hat{\bm{I}}_{\beta} = \tfrac{1}{2}\left(\FrameBasis \otimes\FrameBasis_{\beta} + \FrameBasis_{\beta}\otimes\FrameBasis\right)
\]

\subsubsection{Projections}\label{projections}

The \emph{frame stress} condition is defined by three constraints 
\[
\FrameBasis_{\alpha}\cdot\bm{S}\FrameBasis_{\beta} = 0
\]
\[
\FourTensor{P}_{\mathrm{sec}} = \FrameBasis_{\alpha}\otimes\FrameBasis_{\beta}\otimes\FrameBasis_{\alpha}\otimes\FrameBasis_{\beta}
\quad\text{ and }\quad
\FourTensor{P}_{\mathrm{rod}} = \mathbb{I} - \FourTensor{P}_{\mathrm{sec}}
\]
where 
\[
\mathbb{I} = \FrameBasis_j\otimes\FrameBasis_k \otimes \FrameBasis_j \otimes \FrameBasis_k 
\]
also let 
\[
\FourTensor{P}_{\mathrm{rod}} = \FourTensor{P}_{\mathrm{eul}} + \FourTensor{P}_{\mathrm{cos}}
\]
where 
\[
\FourTensor{P}_{\mathrm{eul}}\triangleq \FrameBasis\otimes\FrameBasis\otimes\FrameBasis\otimes\FrameBasis
\quad\text{ and }\quad 
\FourTensor{P}_{\mathrm{cos}} \triangleq \mathbb{I} - \FourTensor{P}_{\mathrm{eul}}
\]

\subsubsection{Moduli Components}\label{moduli-components}

Introduce a third order tensor
\(\ThreeTensor{P}_{\mathrm{sym}}^{\FrameBasis}\) defined so that 
\[
\ThreeTensor{P}_{\mathrm{sym}}^{\FrameBasis} \bm{\varepsilon} = \FourTensor{P}_{\mathrm{sym}}[\FrameBasis \otimes\bm{\varepsilon}] = \FourTensor{P}_{\mathrm{rod}}\bm{E}
\] 
We can find a representation of
\(\ThreeTensor{P}_{\mathrm{sym}}^{\FrameBasis}\) by appropriately
contracting \(\FourTensor{P}_{\mathrm{sym}}\) with \(\FrameBasis\). 
\[
\ThreeTensor{P}_{\mathrm{sym}}^{\FrameBasis} = \frac{1}{2} \left(\FrameBasis\otimes \FrameBasis_k\otimes\FrameBasis_k + \FrameBasis_k \otimes \FrameBasis\otimes\FrameBasis_k\right)
\] 
we should verify
\(\FourTensor{P}_{\mathrm{sec}}:\ThreeTensor{P}_{\mathrm{sym}}^{\FrameBasis} = \ThreeTensor{0}\).
Then we can write 
\[
\bm{S}\FrameBasis = \ThreeTensor{C}_{\mathrm{rod}}^{\FrameBasis}:\ThreeTensor{P}_{\mathrm{sym}}^{\FrameBasis} \bm{\varepsilon}
\]
